The growing adoption of large language models in production systems makes guardrails, safety classification mechanisms that inspect LLM inputs and outputs, critical components of infrastructures serving users across multiple languages. In local deployment environments, post-training quantization is widely used to reduce memory requirements and increase throughput. However, the impact of quantization on guardrail performance in languages other than English, especially Brazilian Portuguese, remains systematically uninvestigated in the literature. This work evaluates how quantization affects the performance gap between English and Brazilian Portuguese in LLM-based guardrails, operationalized through the $\Delta$F1 metric. Five models are evaluated, Llama Guard 3 8B, Llama Guard 3 1B, ShieldGemma 9B, Nemotron Safety 8B, and Granite Guardian 4.1 8B, across three levels of numerical precision: FP16, GPTQ 8-bit, and GPTQ 4-bit, on a dataset of 16,188 prompts adapted from BeaverTails and translated to PT-BR via Google Translate. The experiment employs Bootstrap with 1,000 iterations for confidence interval estimation and the McNemar test for statistical significance verification of the gap between languages. 
\keywords{guardrails. quantization. LLM safety. Brazilian Portuguese. multilingualism.}